% ================================================================
%  Língua e codificação
% ================================================================
\usepackage[T1]{fontenc}
\usepackage[utf8]{inputenc}
\usepackage[brazil]{babel}

% ================================================================
%  Fonte — Palatino com matemática Pazo
%  (padrão de livros Springer clássicos: Rudin, Lang, etc.)
% ================================================================
\usepackage[osf,sc]{mathpazo}   % old-style figures + versalete
\linespread{1.08}               % Palatino pede entrelinha ligeiramente maior

% ================================================================
%  Matemática
% ================================================================
\usepackage{amsmath, amssymb, amsthm, amsfonts}
\usepackage{mathtools}
\usepackage{calrsfs}

% ================================================================
%  Layout de página
% ================================================================
\usepackage[
	b4paper,
	top    = 3.2cm,
	bottom = 3.8cm,
	inner  = 3.8cm,
	outer  = 2.8cm,
	headsep = 1.2cm
]{geometry}

% ================================================================
%  Tipografia refinada
% ================================================================
\usepackage{microtype}

% Parágrafo clássico: recuo sem espaço extra entre parágrafos
\usepackage{indentfirst}
\setlength{\parindent}{1.3em}
\setlength{\parskip}{0pt plus 0.5pt}

% ================================================================
%  Estilo de capítulo — clássico, centrado, com filete
% ================================================================
\usepackage{titlesec}

% Comando auxiliar: esconde os colchetes de \titlerule[0.4pt]
% do parser dos argumentos do \titleformat — evita conflitos.
\newcommand{\capfilete}{\titlerule[0.4pt]}

% Capítulo
\titleformat{\chapter}[display]
{\normalfont\centering}
{\normalsize\scshape\chaptertitlename\ \thechapter}
{0.6em}
{\capfilete\vspace{0.7em}\huge\bfseries}
[\vspace{0.7em}\capfilete]

\titlespacing*{\chapter}{0pt}{3cm}{2.5cm}

% Seção
\titleformat{\section}
{\normalfont\large\bfseries}
{\thesection}{1em}{}

% Subseção — itálico discreto
\titleformat{\subsection}
{\normalfont\normalsize\bfseries\itshape}
{\thesubsection}{1em}{}

% Parte — majestosa e centrada
\titleformat{\part}[display]
{\normalfont\centering\Huge\scshape}
{\large Parte\ \thepart}
{1em}
{}

% ================================================================
%  Cabeçalho e rodapé — estilo clássico em versalete
% ================================================================
\usepackage{fancyhdr}
\pagestyle{fancy}
\fancyhf{}
\fancyhead[LE]{\small\scshape\nouppercase{\leftmark}}   % capítulo (pág. par)
\fancyhead[RO]{\small\scshape\nouppercase{\rightmark}}  % seção (pág. ímpar)
\fancyfoot[C]{\small\thepage}
\renewcommand{\headrulewidth}{0.3pt}
\renewcommand{\footrulewidth}{0pt}

% Páginas de capítulo: só número centralizado
\fancypagestyle{plain}{%
	\fancyhf{}
	\fancyfoot[C]{\small\thepage}
	\renewcommand{\headrulewidth}{0pt}
}

% ================================================================
%  Figuras, tabelas e floats
% ================================================================
\usepackage{graphicx}
\usepackage{booktabs}
\usepackage{caption}
\usepackage{subcaption}
\usepackage{float}
\captionsetup{
	font        = small,
	labelfont   = {sc},
	format      = plain,
	margin      = 1.2cm,
	labelsep    = period
}

% ================================================================
%  TikZ e algoritmos
% ================================================================
\usepackage{tikz}
\usepackage{tikz-3dplot}
\usepackage[linesnumbered,ruled,vlined,portuguese]{algorithm2e}

% ================================================================
%  Hyperlinks — discretos e clássicos (sem cor)
% ================================================================
\usepackage{bookmark}
\usepackage{hyperref}
\hypersetup{
	colorlinks = true,
	linkcolor  = black,
	citecolor  = black,
	urlcolor   = black,
}

% ================================================================
%  Demonstração e QED
% ================================================================
\renewcommand{\proofname}{Demonstra\c{c}\~{a}o}
\renewcommand{\qedsymbol}{\ensuremath{\blacksquare}}

\makeatletter
\renewenvironment{proof}[1][\proofname]{%
	\par
	\pushQED{\qed}%
	\normalfont\topsep6\p@\@plus6\p@\relax
	\trivlist
	\item[\hskip\labelsep\itshape #1\@addpunct{.}]\ignorespaces
	\setlength{\parindent}{0pt}
}{%
	\popQED\endtrivlist\@endpefalse
}
\makeatother

% ================================================================
%  Operadores e macros matemáticas
% ================================================================
\DeclareMathOperator*{\argmax}{arg\,max}
\newcommand{\mo}[1]{\mathrm{moda}\!\left(#1\right)}
\newcommand{\dm}[1]{\mathrm{dm}\!\left(#1\right)}
\newcommand{\var}[1]{\mathrm{Var}\!\left(#1\right)}
\newcommand{\dep}[1]{\mathrm{dp}\!\left(#1\right)}
\newcommand{\Parts}[1]{\mathcal{P}\!\left(#1\right)}
\newcommand{\borel}{\mathcal{B}(\mathbb{R})}
\newcommand{\real}{\mathbb{R}}
\newcommand{\nat}{\mathbb{N}}
\newcommand{\im}[1]{\mathrm{Im}\!\left(#1\right)}
\newcommand{\norm}[1]{\lVert#1\rVert}
\newcommand{\qs}{\;\;\mathrm{q.s.}\;}
\newcommand{\esssup}{\mathrm{ess}\sup}
\renewcommand{\vec}{\mathbf}
